\subsection*{Quick Example}

The value numbering process assigns a value number to all entities containing a list of values, including for mdspan and ituple. For example,

\begin{lstlisting}[language=croqtile,caption={A shape expression}]
mdspan a : {1, b, c + 3}
\end{lstlisting}

The value numbering process would produce values including:

\begin{lstlisting}[language=croqtile,caption={Its value-number graph}]
   VN #0:  const_1
   VN #1:  b
   VN #2:  c
   VN #3:  const_3
   VN #4:  +:#2:#3
   VN #5:  #0,#1,#4
\end{lstlisting}

A number after "VN \#" identifies an expression's value. It may represent a constant integer, a symbolic value, or a combination of other value numbers. Here, \#0 and \#3 refer to constants; \#1 and \#2 refer to the integer variables \texttt{b} and \texttt{c}. Value \#4 represents the sum of \#2 and \#3, which expands to \texttt{c + 3}. Value \#5 represents the list of values in \texttt{mdspan a}. Indexing \texttt{a(0)} selects \#0, the constant 1.
